% Portada de M2.3: una matriz, su aproximación de bajo rango y los valores singulares.
\documentclass[tikz,border=0pt]{standalone}
\usepackage{resonancia}
\begin{document}
\begin{tikzpicture}
  \clip[rounded corners=6pt] (0,0) rectangle (12,3.6);
  \shade[top color=rAzul!22, bottom color=white] (0,0) rectangle (12,3.6);
  % matriz completa, igual a la suma de capas de rango 1
  \foreach \i in {0,...,5} { \foreach \j in {0,...,5} {
    \pgfmathtruncatemacro{\t}{20+14*abs(mod(\i*2+\j*3,5)-2)}
    \fill[rAzul!\t, draw=white, line width=0.5pt] ({0.5+0.4*\j},{3.0-0.4*\i}) rectangle ({0.5+0.4*(\j+1)},{3.0-0.4*(\i+1)});
  } }
  \node[etiqueta] at (1.7,0.35) {$\mathbf A$};
  \node at (3.5,1.8) {$=$};
  % capas de rango 1: columna por fila
  \foreach \k/\x/\col in {1/4.1/rNaranja, 2/5.9/rVerde, 3/7.7/rVioleta} {
    \fill[\col!55] (\x,2.9) rectangle (\x+0.22,1.0);
    \fill[\col!55] (\x+0.35,2.9) rectangle (\x+1.25,2.68);
    \node[etiqueta, color=\col] at (\x+0.65,0.55) {\small $\sigma_\k\,\mathbf u_\k\mathbf v_\k^{\mathsf T}$};
    \ifnum\k<3 \node at (\x+1.55,1.8) {$+$}; \fi
  }
  % valores singulares
  \draw[eje] (9.6,0.9) -- (11.7,0.9);
  \foreach \x/\h in {9.9/1.9, 10.4/1.0, 10.9/0.45, 11.4/0.1} {
    \fill[rGris!70] (\x-0.15,0.9) rectangle (\x+0.15,{0.9+\h});
  }
  \node[etiqueta, anchor=south] at (10.65,3.0) {\small $\sigma_1\ge\sigma_2\ge\cdots$};
\end{tikzpicture}
\end{document}
